\documentclass[11pt,a4paper]{article}
\usepackage[a4paper,margin=2.5cm]{geometry}
\usepackage{lmodern}
\usepackage[T1]{fontenc}
\usepackage{graphicx}
\setkeys{Gin}{width=0.88\linewidth,height=0.68\textheight,keepaspectratio}
\usepackage{longtable,booktabs,array,etoolbox,calc}
\usepackage{enumitem}
\setlist{topsep=3pt,itemsep=2pt,parsep=0pt}
\usepackage{xcolor}
\usepackage[hidelinks]{hyperref}
\setlength{\emergencystretch}{3em}
\setlength{\parskip}{5pt plus 2pt}
\widowpenalty=10000
\providecommand{\tightlist}{\setlength{\itemsep}{2pt}\setlength{\parskip}{0pt}}
\providecommand{\pandocbounded}[1]{#1}
\begin{document}

\begin{center}
{\LARGE\bfseries Predicting Fair/Poor Self-Rated Health from BRFSS 2021: A Binary Classification Study}\par\vspace{1.0em}
{\normalsize Hassene Bouamama}\par\vspace{0.3em}
{\small Machine Learning final project --- University of Alicante (C2 call)}\par\vspace{0.8em}
{\small\itshape November 2026 --- final draft v1.2 (compiled 2026-11-08)}\par
\end{center}
\vspace{0.6em}

{\small \par}
\vspace{1.2em}
\hrule
\vspace{1.2em}

\tableofcontents
\newpage
\begin{itemize}
\tightlist
\item
  \textbf{Recommended model:} tuned LightGBM --- mean CV ROC-AUC
  \textbf{0.9043} (EXP-014a), nested-CV-confirmed \textbf{0.9043}
  (EXP-015)
\item
  \textbf{Report status: final draft v1.2}
\end{itemize}

\begin{quote}
\textbf{Traceability.} All 42 experiment rows live in the append-only
log \texttt{results/experiments\_log.csv}; every number below cites its
row (e.g.~EXP-014a). Each row was produced under the identical
stratified 10-fold split ---
\texttt{StratifiedKFold(10,\ shuffle=True,\ random\_state=42)},
fold-index SHA \texttt{3b3a17457039}, asserted per run. Descriptive
analyses that fit no model (EDA, Wilcoxon tests, error analysis, SHAP)
carry no log row by design; they cite their committed artifacts, whose
provenance is documented in \texttt{results/p7/oof\_provenance.json}.
\end{quote}

\section{Introduction \& dataset}\label{introduction-dataset}

The Behavioral Risk Factor Surveillance System (BRFSS) is the CDC's
premier system of health-related telephone surveys. It collects
state-level data on U.S. residents' health-related risk behaviors,
chronic conditions, and use of preventive services across all 50 states,
the District of Columbia, and three territories --- more than 400,000
adult interviews per year (course guidelines, \S{}1).

This project uses the BRFSS 2021 LLCP public file: \textbf{438,693
respondent records $\times$ 303 columns}, all numeric-encoded (EXP-000). The
analysis target is \texttt{GENHLTH}, the five-level self-rated health
item.

Self-rated general health is a meaningful prediction target: it is one
of the most widely used summary indicators in survey epidemiology, and
the CDC itself publishes an official derived indicator from it.
Predicting it from behavioral, functional, and socio-demographic factors
is a technically demanding instance of learning from wide, noisy,
sentinel-coded survey data.

The dataset poses four documented challenges, all confirmed empirically
in EXP-000:

\begin{enumerate}
\def\labelenumi{\arabic{enumi}.}
\tightlist
\item
  \textbf{High dimensionality with heavy missingness.} 92 of 303 columns
  exceed 90\% blank+sentinel; 155 exceed 50\%; 3 are completely empty;
  247 contain at least one sentinel code (EXP-000).
\item
  \textbf{Sentinel-coded missingness.} Don't-know/refused values
  (7/77/777, 9/99/999, and 6-digit variants) and codebook-specific
  ``none'' codes (88) are stored as ordinary numbers and must be recoded
  before modeling.
\item
  \textbf{Class imbalance.} Under the chosen binary framing, 16.63\% of
  valid rows are positive (Table 1).
\item
  \textbf{Variables directly related to the target.} Calculated columns
  derived from \texttt{GENHLTH} (notably \texttt{\_RFHLTH}) would leak
  the label; they are identified and blacklisted before modeling
  (EXP-000).
\end{enumerate}

\textbf{Table 1 --- Target distribution (EXP-000).}

{\def\LTcaptype{none} % do not increment counter

\begin{center}\small
\begin{tabular}{@{}
  >{\raggedright\arraybackslash}p{(\linewidth - 6\tabcolsep) * \real{0.2500}}
  >{\raggedright\arraybackslash}p{(\linewidth - 6\tabcolsep) * \real{0.2500}}
  >{\raggedright\arraybackslash}p{(\linewidth - 6\tabcolsep) * \real{0.2500}}
  >{\raggedright\arraybackslash}p{(\linewidth - 6\tabcolsep) * \real{0.2500}}@{}}
\toprule\noalign{}
\begin{minipage}[b]{\linewidth}\raggedright
GENHLTH code
\end{minipage} & \begin{minipage}[b]{\linewidth}\raggedright
Meaning
\end{minipage} & \begin{minipage}[b]{\linewidth}\raggedright
n
\end{minipage} & \begin{minipage}[b]{\linewidth}\raggedright
share of valid
\end{minipage} \\
\midrule\noalign{}
1 & Excellent & 77,741 & 17.8\% \\
2 & Very good & 149,112 & 34.1\% \\
3 & Good & 137,938 & 31.5\% \\
4 & Fair & 54,736 & 12.5\% \\
5 & Poor & 18,005 & 4.1\% \\
7 / 9 / blank & Don't know / Refused / missing & 788 / 369 / 4 & dropped
(1,161 rows) \\
\textbf{y = 1 (Fair/Poor)} & \{4, 5\} & \textbf{72,741} &
\textbf{16.63\%} \\
\textbf{y = 0 (rest)} & \{1, 2, 3\} & \textbf{364,791} &
\textbf{83.37\%} \\
\bottomrule
\end{tabular}
\end{center}

}

After dropping the 1,161 unusable target rows, \textbf{437,532 valid
rows} remain, with a 1 : 5.0 positive:negative ratio (EXP-000).

The work addresses every course learning objective:

\begin{itemize}
\tightlist
\item
  managing and preprocessing a complex real-world clinical dataset;
\item
  framing the native multiclass problem as binary classification
  (ADR-001);
\item
  applying and comparing six classifier families plus a Dummy sanity
  floor;
\item
  robust evaluation --- identical stratified 10-fold CV everywhere,
  Wilcoxon testing, nested CV;
\item
  feature engineering, selection and tuning to maximize ROC-AUC.
\end{itemize}

\section{Task framing --- binary classification per
ADR-001}\label{task-framing-binary-classification-per-adr-001}

\textbf{Target definition (ADR-001).} y = 1 iff \texttt{GENHLTH} $\in$ \{4,
5\} (Fair or Poor); y = 0 iff \texttt{GENHLTH} $\in$ \{1, 2, 3\}. Rows with
\texttt{GENHLTH} $\in$ \{7, 9, blank\} are dropped, never imputed (1,161
rows; EXP-000).

The metric is ROC-AUC, and the validation scheme is stratified 10-fold
CV on identical folds for every model. The fold SHA
\texttt{3b3a17457039} is asserted in all 41 model-experiment rows of
\texttt{results/experiments\_log.csv}; EXP-000, the profiling row, has
no CV folds by design.

\textbf{Why this cut, and why it is not arbitrary.} The CDC itself
derives its official ``adults with fair or poor health'' indicator
\texttt{\_RFHLTH} from \texttt{GENHLTH} with exactly this boundary. Our
binary target therefore mirrors an established operationalization used
in federal reporting rather than a convenience split (ADR-001).

A direct empirical consequence: anything derived from \texttt{GENHLTH}
is target leakage. The leakage probe confirmed this in the strongest
possible terms --- \texttt{\_RFHLTH} agrees with our y on
\textbf{437,532 of 437,532 valid rows (100.0\% agreement, perfectly
diagonal crosstab)} (EXP-000). In the 2021 coding, \texttt{\_RFHLTH} = 2
means Fair/Poor (y = 1); the reversed mapping agrees on zero rows. For
contrast, \texttt{GENHLTH} itself shows Spearman $\rho$ = +0.672 against y
(EXP-000).

The leak columns \texttt{GENHLTH}, \texttt{\_RFHLTH},
\texttt{\_PHYS14D}, \texttt{\_MENT14D} were therefore hard-excluded from
the feature matrix (ADR-004/ADR-005). The matrix is built by
\textbf{whitelist}, so a drop-list typo can never leak the target (\S{}3).

A subtler leakage question was resolved with data rather than decree
(ADR-004). The Tier-3 ``construct-suspicion'' block --- the health-days
items \texttt{PHYSHLTH} ($\rho$ = -0.293 pre-recode), \texttt{MENTHLTH} ($\rho$ =
-0.099), \texttt{POORHLTH} ($\rho$ = -0.104; 48.0\% blank, 50.8\% effective
coverage after sentinel recode) --- is strongly associated with the
target, but none is \emph{derived} from \texttt{GENHLTH}; each measures
a distinct construct (days of poor physical/mental health; days of
activity limitation).

The contrast case is instructive: \texttt{\_RFHLTH}'s 100\% agreement is
what true target derivation looks like; $\rho$ $\approx$ 0.3 is what a legitimately
strong predictor looks like. All three were kept, with
\texttt{POORHLTH}'s final fate left to empirical feature selection
(ADR-004).

\textbf{Alternatives considered and rejected (ADR-001):}

\begin{itemize}
\tightlist
\item
  \emph{Poor (5) vs rest} --- 4.1\% positives (18,005 rows; Table 1):
  harsher imbalance, noisier fold scores, thin error-analysis pool.
  Rejected.
\item
  \emph{Keep the native 5-class problem as the main task} ---
  contradicts the course's explicit binary framing and makes the
  required ROC-AUC/Dummy baseline awkward. Rejected, with the optional
  multiclass extension explicitly declined in ADR-002 (scope: Basic +
  Intermediate + Advanced via boosting libraries, automated search, and
  nested CV instead).
\item
  \emph{Good-or-better vs rest} --- identical partition with flipped
  labels; no distinct experiment.
\end{itemize}

\textbf{Consequences of the framing.} A 16.63\% positive rate gives
stable stratified folds (\textasciitilde7.3k positives per fold) and a
genuinely exercised imbalance question, answered empirically in \S{}6. The
ordinal 3\textbar4 boundary being binarized, meanwhile, turns out to be
the dominant error mode --- analyzed in \S{}9.

\section{Data cleaning \& preprocessing
decisions}\label{data-cleaning-preprocessing-decisions}

Preprocessing is governed by one principle from the Leakage watchlist:
\textbf{the feature matrix is built from an explicit keep-list, never by
blacklisting drops} --- a typo in a drop-list must not leak the target.

Raw data are immutable. The cleaned matrix is checkpointed once to
\texttt{data/interim/clean\_v1.parquet} (438,693 $\times$ 76, sentinels $\rightarrow$ NaN,
engineering recodes applied, no target column) so that no CV fold ever
re-cleans (ADR-005).

\textbf{Whitelist construction: 303 $\rightarrow$ 76 columns via seven rules
(ADR-005).} Every one of the 303 census columns received an explicit
verdict and reason --- no silent drops. Rules applied with precedence R1
\textgreater{} R3 \textgreater{} R2 \textgreater{} R6 \textgreater{} R4
\textgreater{} R5 \textgreater{} keep:

\begin{itemize}
\tightlist
\item
  \textbf{R1 --- leakage blacklist:} \texttt{GENHLTH},
  \texttt{\_RFHLTH}, \texttt{\_PHYS14D}, \texttt{\_MENT14D}; enforced by
  a hard assert in preprocessing.
\item
  \textbf{R2 --- identifiers/admin/weights/telephone-flow:}
  \texttt{SEQNO}, \texttt{\_PSU}, \texttt{\_LLCPWT}, \texttt{\_STSTR},
  MSA/landline/cell admin, \texttt{NUMHHOL3}, \texttt{CPDEMO1B},
  \texttt{MSCODE}.
\item
  \textbf{R3 --- census-constant columns:} 11 constants +
  \texttt{IYEAR}.
\item
  \textbf{R4 --- sparse micro-coverage modules} (\textgreater90\%
  blank+sentinel): per ADR-004 this includes the caregiver block
  CDASSIST/CDHOUSE/CDSOCIAL, whose \textbar $\rho$\textbar{} = 0.28--0.32 on
  \textasciitilde1\% coverage was a selection artifact.
\item
  \textbf{R5 --- coverage gate:} survival requires $\ge$25\% effective
  coverage after sentinel recode; 10--25\% $\rightarrow$ module drop. Borderline
  cases resolved by the rule: \texttt{PREDIAB1} 25.98\% kept,
  \texttt{PDIABTST} 24.86\% dropped, \texttt{SOMALE} 25.23\% the minimum
  kept.
\item
  \textbf{R6 --- duplicate-resolution with keep-raw preference:}
  e.g.~\texttt{\_BMI5} over the height/weight/BMI-flag family;
  \texttt{\_SMOKER3} over
  \texttt{SMOKE100}/\texttt{SMOKDAY2}/\texttt{\_RFSMOK3};
  \texttt{\_AGE80} over age bins; \texttt{\_RACE} over race groupings;
  \texttt{EDUCA}/\texttt{INCOME3}/\texttt{PRIMINSR} over their computed
  bins; \texttt{EXERANY2} over \texttt{\_TOTINDA}; \texttt{HAVARTH5}
  over \texttt{\_DRDXAR3}; alcohol/tobacco/vaccine raws over computed
  indices.
\item
  \textbf{R7 --- recorded judgment calls:} keep \texttt{QSTLANG} drop
  \texttt{QSTVER}; keep \texttt{\_URBSTAT} drop \texttt{\_METSTAT}; keep
  \texttt{FLSHTMY3} as nominal with 777777/999999 $\rightarrow$ NaN; keep
  \texttt{CHOLCHK3} as nominal because its codebook order is
  non-monotone; keep \texttt{CHECKUP1} as ordinal with 8 = ``never'';
  keep \texttt{\_STATE} as a legitimate nominal predictor.
\end{itemize}

\textbf{Table 2 --- Column triage outcome (ADR-005; all 303 columns).}

{\def\LTcaptype{none} % do not increment counter

\begin{center}\small
\begin{tabular}{@{}
  >{\raggedright\arraybackslash}p{(\linewidth - 4\tabcolsep) * \real{0.3333}}
  >{\raggedright\arraybackslash}p{(\linewidth - 4\tabcolsep) * \real{0.3333}}
  >{\raggedright\arraybackslash}p{(\linewidth - 4\tabcolsep) * \real{0.3333}}@{}}
\toprule\noalign{}
\begin{minipage}[b]{\linewidth}\raggedright
Verdict
\end{minipage} & \begin{minipage}[b]{\linewidth}\raggedright
n
\end{minipage} & \begin{minipage}[b]{\linewidth}\raggedright
Examples / notes
\end{minipage} \\
\midrule\noalign{}
KEEP-NOM (one-hot) & 54 & chronic-condition items, demographics, module
survivors, \texttt{\_STATE} (53 FIPS), \texttt{\_RACE},
\texttt{QSTLANG}, \texttt{\_URBSTAT} \\
ENG (engineered numeric) & 15 & PHYSHLTH/MENTHLTH/POORHLTH (88$\rightarrow$0),
CHILDREN, ALCDAY5$\rightarrow$drinking-days/month, AVEDRNK3/DRNK3GE5, MARIJAN1, 6
fruit-veg items$\rightarrow$times/month, \texttt{\_BMI5}$\div$100 \\
KEEP-ORD (ordinal) & 4 & CHECKUP1, EDUCA, INCOME3, \texttt{\_SMOKER3} \\
KEEP-NUM & 3 & JOINPAI2 (0--10), MAXDRNKS, \texttt{\_AGE80} \\
DROP-SPARSE & 86 & \textgreater90\% blank+sentinel modules (hepatitis,
HPV, shingles, caregiver detail, \ldots) \\
DROP-DUP & 65 & computed duplicates of kept raws (\texttt{\_TOTINDA},
\texttt{\_DRDXAR3}, \texttt{\_RFBMI5}, \texttt{\_RFDRHV7}, \ldots) \\
DROP-MODULE & 26 & optional modules under the 25\% coverage gate \\
DROP-ID & 18 & identifiers, dates, weights, strata \\
DROP-MODE & 16 & interview-mode / telephone-flow admin \\
DROP-CONSTANT & 12 & census-constant columns incl.~the 3 empty CFS
items \\
DROP-LEAK & 4 & \texttt{GENHLTH}, \texttt{\_RFHLTH}, \texttt{\_PHYS14D},
\texttt{\_MENT14D} \\
\textbf{Total kept} & \textbf{76} & frozen as matrix v1 \\
\bottomrule
\end{tabular}
\end{center}

}

Two duplicate pairs were \emph{proven} identical rather than assumed:
\texttt{EXERANY2} $\equiv$ \texttt{\_TOTINDA} (identical pairwise n = 436,659
and identical $\rho$ = 0.2465) and \texttt{\_DRDXAR3} $\equiv$ \texttt{HAVARTH5}
(identical n = 434,721, identical $\rho$ = -0.2172) --- both in EXP-000.

\textbf{Sentinel handling: per-column value contracts, not pattern
matching.} The BRFSS codebook's generic 7/9/77/99 scan is a trap, and
the census's naive scan both over- and under-counts (ADR-005):

\begin{itemize}
\tightlist
\item
  \texttt{EMPLOY1} code 7 = ``Retired'' (131,546 rows --- a valid value;
  only 9 is a sentinel);
\item
  \texttt{INCOME3} and \texttt{PRIMINSR} codes 7 and 9 are valid
  brackets/plans (sentinels are 77/99);
\item
  \texttt{IMFVPLA2} 7/9 valid; \texttt{JOINPAI2} 7/9 are valid 0--10
  pain ratings;
\item
  \texttt{\_RACE} 7 = Multiracial; \texttt{\_STATE} 9 = Connecticut;
  \texttt{\_AGE80} 77 = age 77 (no sentinels at all);
\item
  6-digit sentinels 777777/999999 exist in
  \texttt{FLSHTMY3}/\texttt{HIVTSTD3}.
\end{itemize}

Consequently every column has a codebook-exact sentinel set, asserted
post-recode; residual-sentinel scans returned zero violations (ADR-005
sanity gates).

\textbf{Day-count ``none'' codes and the sign-flip evidence.} For
day-count variables, 88 means \emph{zero days}, not missing:
\texttt{PHYSHLTH} carries 88 on 287,796 of 438,693 rows
(\textasciitilde66\%) (EXP-000).

The 88$\rightarrow$0 recode is mandatory (ADR-004), and its importance is
empirically measurable. On raw codes, \texttt{PHYSHLTH}'s Spearman $\rho$ vs
y is \textbf{-0.293}, because the ``none'' sentinel dominates the
numeric range; after the recode it is \textbf{+0.450} --- the strongest
whitelist feature and ultimately the model's \#1 driver (\S{}4, \S{}7, \S{}9)
(EXP-000). Treating sentinels as numbers would have handed every model
an inverted health signal.

\textbf{Missingness after recode and imputation.} Overall matrix
missingness is \textbf{20.54\%} (ADR-005), concentrated in two regimes:
a dense core-survey cluster (median column missingness 5.3\%) and a long
optional-module tail up to 74.8\% (\texttt{SOMALE}; \texttt{PREDIAB1}
74.0\%), with a few intermediate columns in between (\texttt{INCOME3} $\approx$
21\%, \texttt{TRNSGNDR} $\approx$ 44\%, \texttt{POORHLTH} 48.0\%)
(\texttt{results/analysis/eda\_notes.md} \S{}2). Because module assignment
is questionnaire design, missingness is structural.

Imputation is therefore conservative and \textbf{fold-safe} --- always
fit on training folds only inside the CV pipeline: most-frequent
imputation + one-hot for nominal columns, median for
ordinal/numeric/engineered (EXP-001..003 setup).

For the gradient-boosting families a stronger option was verified and
adopted (\S{}6): \textbf{NaN-passthrough}. The one-hot encoder encodes NaN
as its own category (51 of 54 NOM columns carry NaN in matrix v1) and
the boosters route NaN natively at split time --- missingness stays
signal instead of being blurred (EXP-005..007).

All ADR-005 sanity gates passed at matrix build:

\begin{itemize}
\tightlist
\item
  blacklist absent; \texttt{EXERANY2} in / \texttt{\_TOTINDA} out;
  \texttt{HAVARTH5} in / \texttt{\_DRDXAR3} out;
\item
  437,532 valid / 72,741 positives / 16.63\%;
\item
  zero residual sentinels; all kept columns $\ge$25\% effective coverage.
\end{itemize}

The 89 codebook-legal \texttt{\_BMI5} values above 80.0 were kept and
flagged (negligible for trees; 88 among non-missing valid rows in EDA).

\section{Exploratory data analysis}\label{exploratory-data-analysis}

EDA was run on matrix v1 (437,532 valid rows $\times$ 76 columns) by the
data-analyst (\texttt{results/analysis/eda\_notes.md}; figures dpi 150).
It has no log row by design (analysis, not a model run); every figure
was value-re-probed programmatically.

\textbf{Target imbalance (Figure 1).} The binary target is 16.63\%
positive (72,741) vs 83.37\% negative (364,791), a 1 : 5.0 ratio.
Stratified 10-fold CV puts \textasciitilde7.3k positives in every fold,
so AUC estimation is stable; the 83.4\% majority-class accuracy baseline
is trivially high, confirming ROC-AUC as the right metric. The figure
motivated the explicit imbalance experiments of \S{}6 rather than presuming
their answer.

\pandocbounded{\includegraphics[keepaspectratio]{C:/Users/Hassene/Documents/Machine learning final project/figures/eda_target_imbalance.png}}

\emph{\textbf{Figure 1 --- Target distribution.} (a) Raw GENHLTH
classes; (b) binary target y (16.63\% Fair/Poor, 1 : 5.0). Source:
results/analysis/eda\_notes.md \S{}1, data per EXP-000.}

\textbf{Missingness (Figure 2).} Per-column missingness concentrates at
the two ends --- core survey (\textless{} 15\%) vs optional modules
(57--75\%) --- with a handful of intermediates (\texttt{INCOME3} $\approx$ 21\%,
\texttt{TRNSGNDR} $\approx$ 44\%, \texttt{POORHLTH} 48.0\%); median 5.3\%. Four
columns are fully complete (\texttt{SEXVAR}, \texttt{QSTLANG},
\texttt{\_STATE}, \texttt{\_AGE80}); 21 columns are \textgreater{} 25\%
and 19 \textgreater{} 50\% missing.

Two of the strongest single predictors sit in the high-missingness
regime (\texttt{JOINPAI2} 68.1\%, \texttt{POORHLTH} 48.0\%), so their
imputed values deserve extra distrust --- a fact that later shaped the
NaN-passthrough choice (\S{}6) and the interpretation of case B (\S{}9).

\pandocbounded{\includegraphics[keepaspectratio]{C:/Users/Hassene/Documents/Machine learning final project/figures/eda_missingness_matrix.png}}

\emph{\textbf{Figure 2 --- Missingness by column} (matrix v1, valid
rows). Bimodal: core survey vs optional modules. Top-10 highest in
crimson. Source: results/analysis/eda\_notes.md \S{}2.}

\textbf{Univariate associations (Figure 3, Table 3).} Spearman $\rho$ vs y,
pairwise-complete, on the 76 whitelist features (no blacklisted column
can appear by construction).

\textbf{Table 3 --- Top-10 whitelist features by \textbar $\rho$\textbar{}
with y (results/analysis/eda\_notes.md \S{}3).}

{\def\LTcaptype{none} % do not increment counter

\begin{center}\small
\begin{tabular}{@{}
  >{\raggedright\arraybackslash}p{(\linewidth - 6\tabcolsep) * \real{0.2500}}
  >{\raggedright\arraybackslash}p{(\linewidth - 6\tabcolsep) * \real{0.2500}}
  >{\raggedright\arraybackslash}p{(\linewidth - 6\tabcolsep) * \real{0.2500}}
  >{\raggedright\arraybackslash}p{(\linewidth - 6\tabcolsep) * \real{0.2500}}@{}}
\toprule\noalign{}
\begin{minipage}[b]{\linewidth}\raggedright
\#
\end{minipage} & \begin{minipage}[b]{\linewidth}\raggedright
Feature
\end{minipage} & \begin{minipage}[b]{\linewidth}\raggedright
$\rho$
\end{minipage} & \begin{minipage}[b]{\linewidth}\raggedright
Reading
\end{minipage} \\
\midrule\noalign{}
1 & PHYSHLTH & \textbf{+0.450} & more poor-physical-health days $\rightarrow$
Fair/Poor (after 88$\rightarrow$0; was -0.293 on raw codes, EXP-000) \\
2 & DIFFWALK & -0.413 & reverse-coded 1/2 item: difficulty walking $\rightarrow$
Fair/Poor \\
3 & JOINPAI2 & +0.343 & more joint pain (68.1\% missing --- module) \\
4 & DIFFALON & -0.320 & difficulty alone $\rightarrow$ Fair/Poor \\
5 & POORHLTH & +0.307 & more activity-limitation days (48.0\%
missing) \\
6 & LMTJOIN3 & -0.290 & limited by joints \\
7 & DIFFDRES & -0.282 & difficulty dressing \\
8 & INCOME3 & -0.278 & lower income bracket $\rightarrow$ Fair/Poor \\
9 & EMPLOY1 & +0.258 & driven by code 8 ``unable to work'': 2.7\% of y=0
vs 21.6\% of y=1 \\
10 & CHCCOPD3 & -0.252 & COPD (reverse-coded) \\
\bottomrule
\end{tabular}
\end{center}

}

The ceiling \textbar $\rho$\textbar{} $\approx$ 0.45 with no leak-like outlier says
the target is multi-causal --- no single feature separates Fair/Poor, so
gains must come from combining feature blocks, which argues for boosted
ensembles over single shallow models. It also confirms the leakage
controls worked: nothing approaches \texttt{\_RFHLTH}'s 100\% agreement
(\S{}2). At the other end, \texttt{POTATOE1} (-0.0005), \texttt{CRGVEXPT}
(-0.001) and \texttt{\_STATE} (+0.007) are effectively uninformative
marginals.

\pandocbounded{\includegraphics[keepaspectratio]{C:/Users/Hassene/Documents/Machine learning final project/figures/eda_association_top25.png}}

\emph{\textbf{Figure 3 --- Top-25 \textbar Spearman $\rho$\textbar{} with the
binary target} (whitelist features only; blacklisted columns absent by
construction). Source: results/analysis/eda\_notes.md \S{}3.}

\textbf{Correlation structure (Figure 4).} Within the top-40 features
shown in the heatmap, the maximum off-diagonal \textbar $\rho$\textbar{} is
\textbf{0.590} (\texttt{CHOLMED3}--\texttt{TOLDHI3}, cholesterol
medication vs ever-told-high-cholesterol). Only 4 pairs exceed 0.5 ---
also \texttt{EMPLOY1}--\texttt{\_AGE80} 0.562 retirement,
\texttt{BPHIGH6}--\texttt{HOMBPCHK} 0.538,
\texttt{HOMBPCHK}--\texttt{HOMRGCHK} 0.509 --- and none exceeds 0.7.

The full 76$\times$76 matrix peaks higher, at \textbar $\rho$\textbar{} =
\textbf{0.744} (\texttt{AVEDRNK3}--\texttt{MAXDRNKS}, co-asked alcohol
items sharing respondents by survey skip-logic). The redundancy is
contained by the ADR-005 duplicate-pair resolutions and empirically
benign: EXP-012's complement, which holds the alcohol block, adds only
\textasciitilde+0.003 AUC over H1.

Within-block redundancy is moderate (health days $\rho$ +0.234\ldots+0.409;
disability items +0.213\ldots+0.432), so tree importances will show
within-block sharing rather than one dominant member --- a prediction
later confirmed by the gain importances (\S{}7). Three undefined cells
(\texttt{HAVARTH5} $\times$ arthritis-module items) are a skip-logic artifact:
the module is asked only of diagnosed cases, so \texttt{HAVARTH5} is
constant on those pairwise subsets.

\pandocbounded{\includegraphics[keepaspectratio]{C:/Users/Hassene/Documents/Machine learning final project/figures/eda_correlation_heatmap.png}}

\emph{\textbf{Figure 4 --- Spearman correlation heatmap}, top-40
features by \textbar $\rho$\textbar, ordered by block. Max off-diagonal
\textbar $\rho$\textbar{} = 0.590 within this top-40 view; the full 76$\times$76
matrix peaks at 0.744 (\S{}4 text). Source: results/analysis/eda\_notes.md
\S{}4.}

\textbf{Class-conditional profiles (Figure 5).} The positive class
skews:

\begin{itemize}
\tightlist
\item
  older --- \texttt{\_AGE80} mean 53.7 vs 59.9 years;
\item
  lower-income --- median bracket 7 vs 5 of 11;
\item
  functionally limited --- DIFFWALK prevalence 8.9\% vs 47.8\%, a 5.4$\times$
  rate ratio, the cleanest single explainer;
\item
  less active --- EXERANY2 ``any exercise'' 80.2\% vs 51.7\%.
\end{itemize}

Day counts are zero-inflated with heavy right tails (PHYSHLTH mean 1.9
vs 13.8 days; MENTHLTH 3.2 vs 8.7). This epidemiologically coherent
signature is itself a sanity check on the target recode.

Data-quality flags recorded for modeling:

\begin{itemize}
\tightlist
\item
  \texttt{\_AGE80} top-coded (7.85\% exactly 80);
\item
  day-count zero spikes from the 88$\rightarrow$0 recode (PHYSHLTH 65.7\%, MENTHLTH
  61.8\% at exactly 0, shared by both classes);
\item
  deliberate 0.5 ``less-than-monthly'' codes in the fruit-veg items.
\end{itemize}

\pandocbounded{\includegraphics[keepaspectratio]{C:/Users/Hassene/Documents/Machine learning final project/figures/eda_key_features_by_class.png}}

*\textbf{Figure 5 --- Class-conditional distributions} of key features
(y=0 vs y=1): \_AGE80, PHYSHLTH, MENTHLTH, INCOME3, DIFFWALK, \_SMOKER3.
Source: results/analysis/eda\_notes.md \S{}5.*

\textbf{Block-level signal (feeds the \S{}7 hypothesis).} Mean
\textbar $\rho$\textbar{} by EDA grouping: health days \textbf{0.316}
\textgreater{} functional limitation/joint pain \textbf{0.306}
\textgreater{} chronic conditions 0.167 \textgreater{}
socio-demographics 0.144 \textgreater{} access 0.088 \textgreater{}
behaviors 0.082 \textgreater{} diet 0.052 \textgreater{} admin/geo 0.035
(\texttt{results/analysis/eda\_notes.md} \S{}7).

Current-health-status features carry roughly twice the marginal signal
of chronic conditions and four times that of behaviors --- but with
32--48\% less coverage. This asymmetry is the empirical seed of
hypothesis H1.

\section{Baselines}\label{baselines}

All baselines ran on matrix v1 (all-76 features; 437,532 $\times$ 76) under the
canonical folds, with imputation inside the CV pipeline.

\textbf{Dummy floor (EXP-001).}
\texttt{DummyClassifier(strategy=\textquotesingle{}prior\textquotesingle{})}
scored \textbf{ROC-AUC = 0.500000 on every one of the 10 folds} (fold
std 0.000000, runtime 8.3 s) --- exactly what theory predicts for a
constant-score predictor (AUC = $1/2$ under all-ties). This is more than a
formality: it verifies the scorer, fold builder, and target wiring, and
gives every later number a \emph{measured} floor.

\textbf{Logistic Regression (EXP-002a--d).} \texttt{LogisticRegression}
+ \texttt{StandardScaler} over a manual grid C $\in$ \{0.01, 0.1, 1, 10\}
(lbfgs, max\_iter=2000):

\textbf{Table 4 --- LogReg C grid (EXP-002a--d).}

{\def\LTcaptype{none} % do not increment counter

\begin{center}\small
\begin{tabular}{@{}llll@{}}
\toprule\noalign{}
exp\_id & C & mean CV ROC-AUC & fold std \\
\midrule\noalign{}
\textbf{EXP-002a} & \textbf{0.01} & \textbf{0.892588} & 0.001606 \\
EXP-002b & 0.1 & 0.892580 & 0.001607 \\
EXP-002c & 1 & 0.892580 & 0.001607 \\
EXP-002d & 10 & 0.892580 & 0.001607 \\
\bottomrule
\end{tabular}
\end{center}

}

The grid is \emph{flat by three orders of magnitude} --- the top three
settings differ by 8$\times$10\textsuperscript{-}\textsuperscript{6}, far below fold noise (std $\approx$ 0.0016). This is
not under-convergence: a fold-1 refit of the winner converged in 18 of
2000 lbfgs iterations with zero warnings (EXP-002 note).

The interpretation: on \textasciitilde267 one-hot columns with n $\approx$ 394k
training rows per fold, regularization strength is immaterial in this
range. C = 0.01 is kept as the family representative (mild L2 wins by a
whisker and shrinks the wide design). The fold scores of the winner span
only 0.8897--0.8952 --- stable folds, comparisons trustworthy to the
fourth decimal.

\textbf{Decision Tree (EXP-003a--l).} \texttt{DecisionTreeClassifier}
over max\_depth $\in$ \{6, 10, 16, None\} $\times$ min\_samples\_leaf $\in$ \{1, 50,
200\}:

\textbf{Table 5 --- Decision-tree grid (EXP-003a--l), mean CV ROC-AUC.}

{\def\LTcaptype{none} % do not increment counter

\begin{center}\small
\begin{tabular}{@{}
  >{\raggedright\arraybackslash}p{(\linewidth - 14\tabcolsep) * \real{0.1250}}
  >{\raggedright\arraybackslash}p{(\linewidth - 14\tabcolsep) * \real{0.1250}}
  >{\raggedright\arraybackslash}p{(\linewidth - 14\tabcolsep) * \real{0.1250}}
  >{\raggedright\arraybackslash}p{(\linewidth - 14\tabcolsep) * \real{0.1250}}
  >{\raggedright\arraybackslash}p{(\linewidth - 14\tabcolsep) * \real{0.1250}}
  >{\raggedright\arraybackslash}p{(\linewidth - 14\tabcolsep) * \real{0.1250}}
  >{\raggedright\arraybackslash}p{(\linewidth - 14\tabcolsep) * \real{0.1250}}
  >{\raggedright\arraybackslash}p{(\linewidth - 14\tabcolsep) * \real{0.1250}}@{}}
\toprule\noalign{}
\begin{minipage}[b]{\linewidth}\raggedright
exp\_id
\end{minipage} & \begin{minipage}[b]{\linewidth}\raggedright
max\_depth
\end{minipage} & \begin{minipage}[b]{\linewidth}\raggedright
min\_samples\_leaf
\end{minipage} & \begin{minipage}[b]{\linewidth}\raggedright
mean AUC
\end{minipage} & \begin{minipage}[b]{\linewidth}\raggedright
exp\_id
\end{minipage} & \begin{minipage}[b]{\linewidth}\raggedright
max\_depth
\end{minipage} & \begin{minipage}[b]{\linewidth}\raggedright
min\_samples\_leaf
\end{minipage} & \begin{minipage}[b]{\linewidth}\raggedright
mean AUC
\end{minipage} \\
\midrule\noalign{}
EXP-003a & 6 & 1 & 0.848558 & EXP-003g & 16 & 1 & 0.795855 \\
EXP-003b & 6 & 50 & 0.848645 & EXP-003h & 16 & 50 & 0.874830 \\
EXP-003c & 6 & 200 & 0.848666 & \textbf{EXP-003i} & \textbf{16} &
\textbf{200} & \textbf{0.880206} \\
EXP-003d & 10 & 1 & 0.865805 & EXP-003j & None & 50 & 0.864202 \\
EXP-003e & 10 & 50 & 0.871808 & EXP-003k & None & 200 & 0.879195 \\
EXP-003f & 10 & 200 & 0.872352 & EXP-003l & None & 1 & 0.694922 \\
\bottomrule
\end{tabular}
\end{center}

}

Two lessons:

\begin{itemize}
\tightlist
\item
  \textbf{\texttt{min\_samples\_leaf} dominates the grid.} At every
  depth, leaf=200 $\ge$ leaf=50 $\ge$ leaf=1; leaf regularization is what makes
  deep trees usable (with leaf=200, depth 6$\rightarrow$16$\rightarrow$None improves
  0.8487$\rightarrow$0.8802$\rightarrow$0.8792, saturating near 16).
\item
  \textbf{Unconstrained depth overfits catastrophically.} (None, 1) ---
  near scikit-learn defaults --- collapses to \textbf{0.694922}
  (EXP-003l), 0.185 below the winner.
\end{itemize}

The manual grid lifting 0.6949 $\rightarrow$ 0.8802 (EXP-003l $\rightarrow$ EXP-003i) is the
rubric's ``manual hyperparameter adjustment improving ROC-AUC'' evidence
in its cleanest form; the LogReg grid documents the complementary honest
case where a hyperparameter simply does not matter.

The best single tree still loses to regularized linearity by 0.0124 AUC
(0.880206 vs 0.892588) --- expected for one tree vs a linear model on
wide one-hot data, and precisely the gap ensembles exist to close (\S{}6).

\section{Model zoo \& imbalance
handling}\label{model-zoo-imbalance-handling}

\textbf{Model zoo (P4 protocol: one sensible config per family, tuning
deferred to \S{}8).} All runs use all-76 features and the canonical folds.
Random Forest kept the impute preprocessor as a uniform-pipeline choice
(scikit-learn $\ge$ 1.4 can route NaN natively in trees, so this was
convention rather than necessity on the pinned 1.9.1).

The three boosters use \textbf{NaN-passthrough}. It was verified on
sklearn 1.9.1 that
\texttt{OneHotEncoder(handle\_unknown=\textquotesingle{}ignore\textquotesingle{})}
encodes NaN as its own category (51/54 NOM columns carry NaN) and that
HistGB/LightGBM/XGBoost route NaN natively at split time (EXP-005..007).

\textbf{Table 6 --- The zoo: six families + the Dummy floor, seven
configurations, identical folds (fold SHA 3b3a17457039 asserted in every
row).}

{\def\LTcaptype{none} % do not increment counter

\begin{center}\small
\begin{tabular}{@{}
  >{\raggedright\arraybackslash}p{(\linewidth - 8\tabcolsep) * \real{0.2000}}
  >{\raggedright\arraybackslash}p{(\linewidth - 8\tabcolsep) * \real{0.2000}}
  >{\raggedright\arraybackslash}p{(\linewidth - 8\tabcolsep) * \real{0.2000}}
  >{\raggedright\arraybackslash}p{(\linewidth - 8\tabcolsep) * \real{0.2000}}
  >{\raggedright\arraybackslash}p{(\linewidth - 8\tabcolsep) * \real{0.2000}}@{}}
\toprule\noalign{}
\begin{minipage}[b]{\linewidth}\raggedright
exp\_id
\end{minipage} & \begin{minipage}[b]{\linewidth}\raggedright
Model
\end{minipage} & \begin{minipage}[b]{\linewidth}\raggedright
Key config
\end{minipage} & \begin{minipage}[b]{\linewidth}\raggedright
mean CV ROC-AUC $\pm$ fold std
\end{minipage} & \begin{minipage}[b]{\linewidth}\raggedright
runtime
\end{minipage} \\
\midrule\noalign{}
EXP-001 & DummyClassifier & strategy=prior & 0.500000 $\pm$ 0.000000 & 8.3
s \\
EXP-002a & LogisticRegression & C=0.01 + StandardScaler & 0.892588 $\pm$
0.001606 & 13.8 s \\
EXP-003i & DecisionTree & depth 16, leaf 200 & 0.880206 $\pm$ 0.002219 &
32.2 s \\
EXP-004 & RandomForest & 200 trees, leaf 50, sqrt features & 0.893315 $\pm$
0.001530 & 351.5 s \\
EXP-005 & HistGradientBoosting & max\_iter 300, lr 0.06 & 0.903220 $\pm$
0.001748 & 199.9 s \\
EXP-006 & LightGBM & 500 trees, lr 0.05, 63 leaves & 0.903355 $\pm$ 0.001691
& 41.4 s \\
EXP-007 & XGBoost & 500 trees, lr 0.05, depth 6, hist & 0.903357 $\pm$
0.001715 & 138.3 s \\
\bottomrule
\end{tabular}
\end{center}

}

Three findings:

\begin{itemize}
\tightlist
\item
  \textbf{Boosting, not ensembling per se, is the win.} RandomForest
  (bagging) adds only +0.00073 over LogReg --- under half a fold std, a
  statistical tie --- at the worst cost in the zoo (351.5 s, 2.5$\times$
  XGBoost, 8.5$\times$ LightGBM); decorrelated \texttt{sqrt}-feature trees are
  individually too weak on wide one-hot survey data (EXP-004).
\item
  \textbf{The booster trio clears the linear and bagging baselines by a
  real margin:} 0.9032--0.9034 vs 0.8933 at best among
  LogReg/RandomForest --- a +0.011 gain, $\approx$ 6$\times$ the fold std (e.g.~XGB vs
  LogReg: +0.010769 $\approx$ 6.3$\times$ std, EXP-007 note). HistGradientBoosting is
  itself a scikit-learn model; the specialized libraries (LightGBM,
  XGBoost) edge it by $\le$ 0.00014 --- a statistical tie, not a clear.
  Within the trio the spread is $\le$ 0.00014; XGBoost took the family-A
  slot only by the zoo's deterministic tie-break.
\item
  \textbf{LightGBM has the best accuracy/cost ratio} (41.4 s), which
  matters for the compute-heavy search in \S{}8.
\end{itemize}

The NaN-passthrough result is itself a finding: the boosters ran with
\emph{zero imputation}, keeping module missingness as signal --- a core
reason this family suits sparse survey data (EXP-005..007).

\textbf{Imbalance handling (EXP-008/009/010).} Three techniques $\times$ two
families (XGBoost as family A, LogReg as family B), all resampling
\textbf{inside} the training folds via imblearn pipelines:

\textbf{Table 7 --- Imbalance strategies: $\Delta$AUC vs the same model with
default weights on the same folds.}

{\def\LTcaptype{none} % do not increment counter

\begin{center}\small
\begin{tabular}{@{}
  >{\raggedright\arraybackslash}p{(\linewidth - 4\tabcolsep) * \real{0.3333}}
  >{\raggedright\arraybackslash}p{(\linewidth - 4\tabcolsep) * \real{0.3333}}
  >{\raggedright\arraybackslash}p{(\linewidth - 4\tabcolsep) * \real{0.3333}}@{}}
\toprule\noalign{}
\begin{minipage}[b]{\linewidth}\raggedright
Strategy
\end{minipage} & \begin{minipage}[b]{\linewidth}\raggedright
Family A (XGBoost, baseline EXP-007 = 0.903357)
\end{minipage} & \begin{minipage}[b]{\linewidth}\raggedright
Family B (LogReg, baseline EXP-002a = 0.892588)
\end{minipage} \\
\midrule\noalign{}
none (default) & 0.903357 (EXP-007) & 0.892588 (EXP-002a) \\
Class weighting & EXP-008a: 0.903261 ($\Delta$ -0.000096) & EXP-008b: 0.893035
($\Delta$ +0.000447) \\
SMOTE(0.5) inside folds & EXP-009a: 0.899803 ($\Delta$ -0.003554) & EXP-009b:
0.888419 ($\Delta$ -0.004169) \\
RandomUnderSampler(0.5) & EXP-010a: 0.903096 ($\Delta$ -0.000261) & EXP-010b:
0.892765 ($\Delta$ +0.000177) \\
\bottomrule
\end{tabular}
\end{center}

}

Details worth recording:

\begin{itemize}
\tightlist
\item
  \texttt{scale\_pos\_weight} = neg/pos = 364,791/72,741 $\approx$ 5.0149
  (stratified folds preserve the ratio);
\item
  SMOTE used sampling\_strategy 0.5 because a 1:1 rebalance would
  synthesize 262,846 rows per fold (393,778 $\rightarrow$ 656,624) vs 98,690 at 0.5
  (EXP-009 note);
\item
  RUS drops the per-fold majority class 328,312 $\rightarrow$ 130,932 rows (393,778
  $\rightarrow$ 196,398) (EXP-010 note).
\end{itemize}

\textbf{Verdict: \texttt{imbalance\_strategy\ =\ none} for the AUC
objective.} Class weighting and undersampling move AUC by at most
0.00045 ($\approx$ 0.3$\times$ fold std, both directions). SMOTE is the \emph{worst}
strategy on both families --- the only P4 imbalance configuration where
a booster loses \textasciitilde2$\times$ fold std (-0.00355) and the only
configuration of the entire imbalance study to fall below the plain
LogReg baseline it was meant to improve (EXP-009b, -0.00417).

The reason is the metric, not the implementations: \textbf{ROC-AUC is
threshold-free and rank-based.} Reweighting and resampling reshape the
class prior and move the implicit operating point, but they do not
improve the \emph{ranking} of positives above negatives --- which is all
AUC measures. SMOTE additionally blurs the boundary with synthetic
interpolated minority points (EXP-008..010 verdicts).

One honest caveat is documented: EXP-009a was forced onto the impute
preprocessor (SMOTE rejects NaN), so its $\Delta$ conflates strategy with
preprocessing --- but EXP-009b, which kept its baseline's preprocessing,
shows the damage is real for SMOTE itself.

A side-finding adopted later: RUS(0.5) is AUC-neutral while 1.9$\times$ faster
(73.0 vs 138.3 s for identical XGBoost; EXP-010a) --- a legitimate
acceleration lever for search phases, with finals always re-validated on
full folds.

\section{Feature selection \&
hypothesis}\label{feature-selection-hypothesis}

P5 asked three questions with one protocol. Every selection, importance
and projection is computed \textbf{inside} the 10 training folds only
(never on full data or test rows), then scored on the canonical folds;
baseline rows were never re-run.

\textbf{Hypothesis H1 (stated before the subset experiments,
coordinator-approved; EXP-012):}

\begin{quote}
*The model's predictive power concentrates in five blocks: B1 functional
limitation/disability (DIFFWALK, DIFFALON, DIFFDRES, DECIDE, DEAF,
BLIND); B2 health days (PHYSHLTH, MENTHLTH, POORHLTH); B3 joint
pain/arthritis (JOINPAI2, LMTJOIN3, HAVARTH5, ARTHDIS2, ARTHEXER); B4
chronic conditions (DIABETE4, BPHIGH6, TOLDHI3, CHOLMED3, CVDINFR4,
CVDCRHD4, CVDSTRK3, ASTHMA3, ASTHNOW, CHCSCNCR, CHCOCNCR, CHCCOPD3,
ADDEPEV3, CHCKDNY2); B5 socio-demographic and lifestyle factors
(\_AGE80, INCOME3, EDUCA, EMPLOY1, MARITAL, \_RACE, SEXVAR, and the
lifestyle members \_BMI5, \_SMOKER3). A compact subset of these
\textasciitilde35 columns should match all-76 within fold noise; the
complement (diet, screening, state, remaining module vars) should add
little \emph{on top of the blocks} --- its standalone signal may be
substantial.*
\end{quote}

The hypothesis is grounded in EDA (\S{}4 block ranking: health days and
functional limitation carry \textasciitilde2$\times$ the marginal signal of
chronic conditions) rather than in post-hoc rationalization.

\textbf{Basic selection: SelectKBest (mutual information) --- an honest
negative (EXP-011a).} \texttt{SelectKBest(mutual\_info\_classif)} was
fit inside each training fold on the encoded matrix; k = 50/100/200 are
top-k prefixes of one per-fold MI ranking.

Every k is \emph{below} all-76: k=50 $\rightarrow$ 0.889178, k=100 $\rightarrow$ 0.899556, k=200
$\rightarrow$ 0.899670 (EXP-011a-50/100/200) vs 0.903355 for all-76 (EXP-006/011b).
The best filter (k=200) still costs -0.0037.

The mechanism: MI is univariate, so it drops correlated-but-useful
columns that boosted trees exploit for free --- redundancy costs a tree
nothing. Selection is also unstable at small k: only 48 of 76 columns
are selected in all 10 folds at k=50 (EXP-011 artifact).

\textbf{Advanced selection: fold-internal LightGBM gain importances
(EXP-011b).} Gain shares averaged over the 10 training folds (one-hot
children summed back to the original 76 columns), in the run whose fold
AUCs replicate EXP-006 to 6 decimals (max \textbar diff\textbar{} =
3.7$\times$10\textsuperscript{-}\textsuperscript{7}):

\textbf{Table 8 --- LightGBM gain top-10 (EXP-011b; Figure 6).}

{\def\LTcaptype{none} % do not increment counter

\begin{center}\small
\begin{tabular}{@{}llll@{}}
\toprule\noalign{}
\# & Column & Block & Share of total gain \\
\midrule\noalign{}
1 & PHYSHLTH & B2 & \textbf{51.69\%} \\
2 & DIFFWALK & B1 & 7.99\% \\
3 & DIABETE4 & B4 & 3.31\% \\
4 & INCOME3 & B5 & 2.55\% \\
5 & EDUCA & B5 & 2.53\% \\
6 & \_BMI5 & B5 & 2.44\% \\
7 & POORHLTH & B2 & 1.80\% \\
8 & CHCCOPD3 & B4 & 1.71\% \\
9 & EXERANY2 & other & 1.56\% \\
10 & MENTHLTH & B2 & 1.53\% \\
\bottomrule
\end{tabular}
\end{center}

}

Gain is extraordinarily concentrated --- one column carries half of it
--- and sits overwhelmingly inside the H1 blocks. Block shares: B2
\textbf{55.0\%}, B5 11.7\%, B4 10.6\%, B1 10.1\%, other 11.4\%, B3 1.2\%
--- B3 (joint pain/arthritis) was overweighted by the hypothesis: its
standalone association (\S{}4) largely did not survive competition with
B2's health days. 9 of the top-10 gain columns are in H1-36
(EXP-011/012).

The two importance notions disagree at the top, exactly as different
methods should: MI $\cap$ gain overlap is 3/10 at top-10, 13/30 at top-30
(EXP-011) --- MI ranks coverage-rich items higher, gain ranks
split-useful items higher.

\pandocbounded{\includegraphics[keepaspectratio]{C:/Users/Hassene/Documents/Machine learning final project/figures/p5_importance_top30.png}}

\emph{\textbf{Figure 6 --- LightGBM gain importances, top 30 original
columns} (fold-internal, one-hot children summed back; PHYSHLTH
51.69\%). Source: EXP-011b.}

\textbf{Testing H1: driver-block subset vs all-76 vs complement
(EXP-012a--d).} H1 names 37 columns; 36 exist in matrix v1
(\texttt{ASTHNOW} is absent from the 2021 whitelist --- \texttt{ASTHMA3}
is the sole asthma item). The complement is the remaining 40 columns:
diet/fruit-veg, screening/immunization, alcohol, sexual-orientation and
caregiver modules, admin/geo, access/utilization. Both booster configs
run unchanged:

\textbf{Table 9 --- H1-subset vs all-76 vs complement (EXP-012;
baselines EXP-006/007).}

{\def\LTcaptype{none} % do not increment counter

\begin{center}\small
\begin{tabular}{@{}
  >{\raggedright\arraybackslash}p{(\linewidth - 8\tabcolsep) * \real{0.2000}}
  >{\raggedright\arraybackslash}p{(\linewidth - 8\tabcolsep) * \real{0.2000}}
  >{\raggedright\arraybackslash}p{(\linewidth - 8\tabcolsep) * \real{0.2000}}
  >{\raggedright\arraybackslash}p{(\linewidth - 8\tabcolsep) * \real{0.2000}}
  >{\raggedright\arraybackslash}p{(\linewidth - 8\tabcolsep) * \real{0.2000}}@{}}
\toprule\noalign{}
\begin{minipage}[b]{\linewidth}\raggedright
Feature set
\end{minipage} & \begin{minipage}[b]{\linewidth}\raggedright
Model
\end{minipage} & \begin{minipage}[b]{\linewidth}\raggedright
mean CV ROC-AUC
\end{minipage} & \begin{minipage}[b]{\linewidth}\raggedright
$\Delta$ vs all-76
\end{minipage} & \begin{minipage}[b]{\linewidth}\raggedright
folds subset $\ge$ full
\end{minipage} \\
\midrule\noalign{}
H1-subset (36 cols) & LGBM (EXP-012a) & 0.900051 & -0.00330 & 0/10 \\
H1-subset (36 cols) & XGB (EXP-012b) & 0.900316 & -0.00304 & 0/10 \\
Complement (40 cols) & LGBM (EXP-012c) & 0.811844 & -0.09151 & 0/10 \\
Complement (40 cols) & XGB (EXP-012d) & 0.811881 & -0.09148 & 0/10 \\
all-76 (baseline) & LGBM (EXP-006) & 0.903355 & --- & --- \\
all-76 (baseline) & XGB (EXP-007) & 0.903357 & --- & --- \\
\bottomrule
\end{tabular}
\end{center}

}

\textbf{H1 is nearly confirmed but not strictly within fold noise ---
and the report says so.} The 36-column subset retains \textbf{99.2\% of
the all-76 lift over 0.5} ((0.900051-0.5)/(0.903355-0.5)), on both
families, with paired per-fold $\Delta$ = -0.00330 $\pm$ 0.00050 (EXP-012). The
deficit is tiny but \emph{consistent}: the subset wins \textbf{0 of 10
folds} on both families, and it is not noise (\textasciitilde1.9$\sigma$
against a one-fold-std criterion of 0.001715).

The complement is dramatically weaker standalone (77.3\% of the full
lift retained) but carries real signal, largely redundant with H1
through correlates like EXERANY2 and screening items, plus
module-assignment missingness artifacts. Its \emph{incremental} value on
top of H1 is exactly the \textasciitilde+0.003 the subset gives up
(all-76 = H1 $\cup$ complement on the same folds).

\textbf{Decision: keep all-76 features for the final model; use H1
blocks to drive interpretation} (\S{}9 error analysis and XAI are organized
around them). The rubric's selected-vs-all comparison is thus answered
in both directions with paired folds.

\textbf{Dimensionality reduction: PCA --- a second honest negative
(EXP-013).} Impute $\rightarrow$ StandardScaler $\rightarrow$ PCA (fit inside folds) $\rightarrow$ LightGBM,
at n $\in$ \{50, 100, 200\} components:

\textbf{Table 10 --- PCA vs raw 76 columns (EXP-013; cumulative variance
is the fold mean).}

{\def\LTcaptype{none} % do not increment counter

\begin{center}\small
\begin{tabular}{@{}
  >{\raggedright\arraybackslash}p{(\linewidth - 6\tabcolsep) * \real{0.2500}}
  >{\raggedright\arraybackslash}p{(\linewidth - 6\tabcolsep) * \real{0.2500}}
  >{\raggedright\arraybackslash}p{(\linewidth - 6\tabcolsep) * \real{0.2500}}
  >{\raggedright\arraybackslash}p{(\linewidth - 6\tabcolsep) * \real{0.2500}}@{}}
\toprule\noalign{}
\begin{minipage}[b]{\linewidth}\raggedright
n\_components
\end{minipage} & \begin{minipage}[b]{\linewidth}\raggedright
cum. explained variance
\end{minipage} & \begin{minipage}[b]{\linewidth}\raggedright
mean CV ROC-AUC
\end{minipage} & \begin{minipage}[b]{\linewidth}\raggedright
$\Delta$ vs EXP-006
\end{minipage} \\
\midrule\noalign{}
50 & 43.7\% & 0.881126 & -0.02223 \\
100 & 63.7\% & 0.886203 & -0.01715 \\
200 & 97.8\% & 0.888800 & -0.01456 \\
\bottomrule
\end{tabular}
\end{center}

}

PCA costs \textbf{0.0146--0.0222 AUC} even at the best setting (n=200).
Two reasons: unsupervised linear compression smears the one-hot
structure (the 53-level \texttt{\_STATE} block in particular) into
shared components; and PCA forces imputation, forfeiting the boosters'
native NaN routing that \S{}6 showed is part of their advantage.

An interesting side-finding tempers the negative: 50 components carry
only 43.7\% of variance yet already reach 0.8811 --- the \emph{signal}
is low-dimensional even though the optimal \emph{representation} is not
(Figure 7). Raw columns are kept.

\pandocbounded{\includegraphics[keepaspectratio]{C:/Users/Hassene/Documents/Machine learning final project/figures/p5_pca_cumvar.png}}

\emph{\textbf{Figure 7 --- PCA cumulative explained variance vs number
of components} (fold-internal fits). Source: EXP-013.}

\section{Hyperparameter tuning \& statistical
comparison}\label{hyperparameter-tuning-statistical-comparison}

\textbf{Automated search (EXP-014a/b).}
\texttt{RandomizedSearchCV(n\_iter=15,\ scoring=\textquotesingle{}roc\_auc\textquotesingle{},\ random\_state=42)}
over the two top boosters, fit on the full X,y with inner
\texttt{StratifiedKFold(5,\ shuffle=True,\ random\_state=42)} for
\emph{selection only}. The winners were then scored on the canonical 10
folds via the standard harness (fold SHA asserted).

Search spaces: n\_estimators \textasciitilde{} randint{[}300,600{]},
learning\_rate \textasciitilde{} logU{[}0.02,0.15{]}, subsample
\textasciitilde{} U{[}0.6,1.0{]}, colsample\_bytree \textasciitilde{}
U{[}0.5,1.0{]}, reg\_lambda \textasciitilde{} logU{[}0.1,10{]}, plus
num\_leaves \textasciitilde{} randint{[}15,127{]} and
min\_child\_samples \textasciitilde{} randint{[}5,100{]} (LGBM), or
max\_depth \textasciitilde{} randint{[}3,10{]} and min\_child\_weight
\textasciitilde{} randint{[}1,20{]} (XGB).

Compute-budget levers were documented, not silent. A timing probe plus
P4-calibrated projections put the spec'd n\_iter=20 / n\_estimators$\le$800
worst case at \textasciitilde110 min, over the 90-min budget, so n\_iter
was capped at 15 and the tree range at {[}300,600{]} (untuned 500-tree
configs stay bracketed). Actual P6 wall time was 53.3 min (EXP-014
note).

\textbf{Table 11 --- Search winners (inner 5-fold AUC).}

{\def\LTcaptype{none} % do not increment counter

\begin{center}\small
\begin{tabular}{@{}
  >{\raggedright\arraybackslash}p{(\linewidth - 4\tabcolsep) * \real{0.3333}}
  >{\raggedright\arraybackslash}p{(\linewidth - 4\tabcolsep) * \real{0.3333}}
  >{\raggedright\arraybackslash}p{(\linewidth - 4\tabcolsep) * \real{0.3333}}@{}}
\toprule\noalign{}
\begin{minipage}[b]{\linewidth}\raggedright
Family
\end{minipage} & \begin{minipage}[b]{\linewidth}\raggedright
Best inner AUC $\pm$ std
\end{minipage} & \begin{minipage}[b]{\linewidth}\raggedright
Winning parameters (\texttt{results/p6/best\_params.json})
\end{minipage} \\
\midrule\noalign{}
LightGBM & 0.904222 $\pm$ 0.001051 & n\_estimators=512,
learning\_rate=0.0219, num\_leaves=94, min\_child\_samples=44,
subsample=0.8159, colsample\_bytree=0.5980, reg\_lambda=3.123 \\
XGBoost & 0.903679 $\pm$ 0.001059 & n\_estimators=451,
learning\_rate=0.0505, max\_depth=7, min\_child\_weight=4,
subsample=0.6226, colsample\_bytree=0.5500, reg\_lambda=2.003 \\
\bottomrule
\end{tabular}
\end{center}

}

\textbf{Table 12 --- Tuned finals on the canonical folds (EXP-014 vs
untuned P4 baselines).}

{\def\LTcaptype{none} % do not increment counter

\begin{center}\small
\begin{tabular}{@{}
  >{\raggedright\arraybackslash}p{(\linewidth - 8\tabcolsep) * \real{0.2000}}
  >{\raggedright\arraybackslash}p{(\linewidth - 8\tabcolsep) * \real{0.2000}}
  >{\raggedright\arraybackslash}p{(\linewidth - 8\tabcolsep) * \real{0.2000}}
  >{\raggedright\arraybackslash}p{(\linewidth - 8\tabcolsep) * \real{0.2000}}
  >{\raggedright\arraybackslash}p{(\linewidth - 8\tabcolsep) * \real{0.2000}}@{}}
\toprule\noalign{}
\begin{minipage}[b]{\linewidth}\raggedright
Family
\end{minipage} & \begin{minipage}[b]{\linewidth}\raggedright
Untuned
\end{minipage} & \begin{minipage}[b]{\linewidth}\raggedright
Tuned
\end{minipage} & \begin{minipage}[b]{\linewidth}\raggedright
$\Delta$ AUC
\end{minipage} & \begin{minipage}[b]{\linewidth}\raggedright
Paired folds better
\end{minipage} \\
\midrule\noalign{}
LightGBM & EXP-006: 0.903355 $\pm$ 0.001691 & \textbf{EXP-014a: 0.904321 $\pm$
0.001675} & +0.000966 & \textbf{10/10} \\
XGBoost & EXP-007: 0.903357 $\pm$ 0.001715 & EXP-014b: 0.903843 $\pm$ 0.001667 &
+0.000486 & 9/10 (only fold 2 favours untuned) \\
\bottomrule
\end{tabular}
\end{center}

}

\pandocbounded{\includegraphics[keepaspectratio]{C:/Users/Hassene/Documents/Machine learning final project/figures/p6_tuning_results.png}}

\emph{\textbf{Figure 8 --- P6 tuning results} (search candidates,
tuned-vs-untuned fold scores). Source: EXP-014/015.}

\textbf{Nested cross-validation: the unbiased estimate (EXP-015).}
Because the search saw the full X,y, the Part-1 numbers carry selection
bias by construction --- exactly what nested CV prices in. Outer loop =
the canonical 10 folds (SHA-asserted); inner =
\texttt{RandomizedSearchCV(n\_iter=12,\ StratifiedKFold(3,\ rs=42),\ roc\_auc)}
refit per outer fold on the outer \emph{training} fold only. Result:
\textbf{0.904303 $\pm$ 0.001700}.

The optimism of the non-nested estimate is \textbf{+0.000018} (0.904321
- 0.904303) --- negligible: with 437,532 rows and a narrow, regularized
space, inner-CV selection is nearly deterministic across resamples. Note
the two procedures differ by design: the inner search was deliberately
smaller than Part-1's (12 candidates $\times$ 3 inner folds vs 15 $\times$ 5), so
``7/10 outer folds re-selected the Part-1 winner'' compares winners of
different-budget searches --- and stability held anyway. Against the
untuned EXP-006 baseline the nested tuned gain survives at +0.000949.

The report therefore quotes both numbers with distinct meanings:
\textbf{0.904321} (EXP-014a) is the canonical-folds tuned result,
comparable to every other log row; \textbf{0.904303} (EXP-015) is the
selection-bias-free estimate.

\textbf{Statistical comparison (Wilcoxon signed-rank, exact, on the 10
paired canonical-fold AUCs; analysis note, no log row).} Pairing is
exact because the folds are identical by construction in every feeding
row.

\textbf{Table 13 --- Wilcoxon signed-rank tests (two-sided,
\texttt{zero\_method=\textquotesingle{}wilcox\textquotesingle{}},
exact).}

{\def\LTcaptype{none} % do not increment counter

\begin{center}\small
\begin{tabular}{@{}
  >{\raggedright\arraybackslash}p{(\linewidth - 14\tabcolsep) * \real{0.1250}}
  >{\raggedright\arraybackslash}p{(\linewidth - 14\tabcolsep) * \real{0.1250}}
  >{\raggedright\arraybackslash}p{(\linewidth - 14\tabcolsep) * \real{0.1250}}
  >{\raggedright\arraybackslash}p{(\linewidth - 14\tabcolsep) * \real{0.1250}}
  >{\raggedright\arraybackslash}p{(\linewidth - 14\tabcolsep) * \real{0.1250}}
  >{\raggedright\arraybackslash}p{(\linewidth - 14\tabcolsep) * \real{0.1250}}
  >{\raggedright\arraybackslash}p{(\linewidth - 14\tabcolsep) * \real{0.1250}}
  >{\raggedright\arraybackslash}p{(\linewidth - 14\tabcolsep) * \real{0.1250}}@{}}
\toprule\noalign{}
\begin{minipage}[b]{\linewidth}\raggedright
Test
\end{minipage} & \begin{minipage}[b]{\linewidth}\raggedright
Comparison
\end{minipage} & \begin{minipage}[b]{\linewidth}\raggedright
Mean AUC a vs b
\end{minipage} & \begin{minipage}[b]{\linewidth}\raggedright
Median paired $\Delta$
\end{minipage} & \begin{minipage}[b]{\linewidth}\raggedright
Folds a\textgreater b
\end{minipage} & \begin{minipage}[b]{\linewidth}\raggedright
W
\end{minipage} & \begin{minipage}[b]{\linewidth}\raggedright
p
\end{minipage} & \begin{minipage}[b]{\linewidth}\raggedright
Verdict
\end{minipage} \\
\midrule\noalign{}
T1 & tuned LGBM vs tuned XGB & 0.904321 vs 0.903843 & +0.000471 & 10/10
& 0 & \textbf{0.001953} & EXP-014a \textgreater{} EXP-014b \\
T2 & tuned LGBM vs untuned self & 0.904321 vs 0.903355 & +0.000991 &
10/10 & 0 & \textbf{0.001953} & EXP-014a \textgreater{} EXP-006 \\
T3 & tuned LGBM vs LogReg & 0.904321 vs 0.892588 & +0.011768 & 10/10 & 0
& \textbf{0.001953} & EXP-014a \textgreater{} EXP-002a \\
\bottomrule
\end{tabular}
\end{center}

}

All 10 paired differences share the same sign in every test, so W = 0
and p = 0.001953 = 2$\cdot$2\textsuperscript{-}\textsuperscript{1}\textsuperscript{0} --- the smallest two-sided p attainable at n =
10 (the ``n=10 floor''): no 10-fold paired comparison could be more
significant. Three-for-three significant at $\alpha$ = 0.05.

\textbf{The headline is the effect-size contrast: choosing the right
model family (T3, +0.012) buys an order of magnitude more than tuning it
(T2, +0.001).} Both gains are real, but the budget priority is
unambiguous. Tuned LightGBM is confirmed as the final model.

\section{Error analysis \& XAI}\label{error-analysis-xai}

Error analysis and XAI derive deterministically from the already-logged
configurations EXP-014a (tuned LGBM), EXP-014b (tuned XGB) and EXP-002a
(LogReg). Out-of-fold (OOF) predictions come from re-fitting exactly
those pipelines on the canonical folds --- per-fold OOF AUCs re-verified
against the logged fold scores to $\le$ 4.8$\times$10\textsuperscript{-}\textsuperscript{7} for all three models
(\texttt{results/p7/oof\_provenance.json}). SHAP values come from a
single full-data refit of the EXP-014a config used as the
\emph{final-model proxy} --- a different stated object, never mixed with
OOF probabilities. The authoritative interpretation note is
\texttt{results/analysis/error\_analysis.md}.

\textbf{Error taxonomy at threshold 0.5 (tuned LGBM, OOF, n = 437,532).}
TP/FP/TN/FN = 38,957 / 14,340 / 350,451 / 33,784 $\rightarrow$ \textbf{FPR 3.93\%,
FNR 46.44\%} (TPR 0.5356, TNR 0.9607), overall error rate 11.00\%
(48,124 errors) at prevalence 16.63\%
(\texttt{results/p7/error\_taxonomy.csv}).

\textbf{Table 14 --- Errors by raw GENHLTH class: the 3\textbar4
boundary story.}

{\def\LTcaptype{none} % do not increment counter

\begin{center}\small
\begin{tabular}{@{}lllll@{}}
\toprule\noalign{}
GENHLTH class & n & Error type & Count & Rate within class \\
\midrule\noalign{}
1 Excellent & 77,741 & FP & 709 & 0.0091 \\
2 Very good & 149,112 & FP & 2,430 & 0.0163 \\
3 Good & 137,938 & FP & 11,201 & \textbf{0.0812} \\
4 Fair & 54,736 & FN & 31,248 & \textbf{0.5709} \\
5 Poor & 18,005 & FN & 2,536 & 0.1408 \\
\bottomrule
\end{tabular}
\end{center}

}

The FP rate climbs monotonically toward the boundary (0.9\% $\rightarrow$ 1.6\% $\rightarrow$
8.1\%); the FN rate falls as health worsens (57.1\% $\rightarrow$ 14.1\%; class-4
TPR is 0.429 vs class-5 0.859). \textbf{42,449 of 48,124 errors (88.2\%)
sit exactly at the 3\textbar4 boundary} (class-3 FPs + class-4 FNs).

Moreover, 41,091 rows (9.4\% of all) are misclassified by \emph{all
three} architecturally different families simultaneously, while 86.9\%
are wrong in none --- the errors are structured, not stochastic: if they
were noise they would spread evenly across classes and models.

This is the signature of \textbf{boundary ambiguity in the binarized
ordinal label}, not model capacity. ``Fair'' is a heterogeneous
catch-all adjacent to ``Good'', and the self-rating depends on the
respondent's internal reference frame, which the 76 columns cannot
observe.

\textbf{Why FNR 46\% vs FPR 4\% is expected geometry, not a defect.} The
model was selected purely on ROC-AUC with no weighting or resampling
(per \S{}6) at a 1 : 5 prior. AUC optimizes ranking across all thresholds
and says nothing about where the probability mass sits relative to an
arbitrary 0.5 cut; with an abundant negative class the natural operating
point is cheap-on-FPs, expensive-on-FNs.

This is also why resampling did not help AUC in \S{}6 --- and it means 0.5
is simply the wrong knob \emph{if} the model were used for screening (it
is not; course project). The threshold menu from the OOF sweep
(\texttt{results/analysis/error\_analysis.md} \S{}5):

\textbf{Table 15 --- Threshold menu (tuned LGBM OOF; smallest $\tau$
attaining each TPR target).}

{\def\LTcaptype{none} % do not increment counter

\begin{center}\small
\begin{tabular}{@{}lllll@{}}
\toprule\noalign{}
$\tau$ & TPR & FPR & Precision & Share of population flagged \\
\midrule\noalign{}
0.5000 & 0.536 & 0.039 & 0.731 & 12.2\% \\
0.4172 & 0.600 & 0.055 & 0.683 & 14.6\% \\
0.2886 & 0.700 & 0.093 & 0.600 & 19.4\% \\
0.1771 & 0.800 & 0.158 & 0.502 & 26.5\% \\
0.0878 & 0.900 & 0.283 & 0.388 & 38.6\% \\
\bottomrule
\end{tabular}
\end{center}

}

A recall-oriented triage policy would read this trade-off explicitly ---
e.g.~$\tau$ = 0.2886 buys TPR 0.70 at FPR 0.093 (precision 0.60) --- with the
right $\tau$ depending on a use-case cost ratio the course project does not
have.

\textbf{Three archetype cases} (each wrong in 3/3 model families; full
vectors and z-scores in \texttt{results/p7/selected\_cases.json};
profiles in Figure 9, local SHAP decompositions in Figures 10--12):

\begin{itemize}
\tightlist
\item
  \textbf{Case A --- most confident false positive} (row 307041, fold
  4): y = 0 (GENHLTH = 3 ``Good'') yet P(y=1) = 0.9595 / 0.9815 / 0.9930
  (LGBM/XGB/LR). The respondent reports PHYSHLTH = 30 days (z +3.14;
  SHAP \textbf{+1.94}, the largest single contribution in any case),
  POORHLTH = 30 days, difficulty walking, six named chronic conditions
  (diabetes, heart attack, coronary disease, stroke, kidney disease,
  arthritis; the raw vector additionally reports hypertension and high
  cholesterol), BMI 51.45 (z +3.50), lowest income bracket, retired.
  Net: base -2.6095 + $\Sigma$SHAP +5.64 = log-odds +3.03 (proxy P = 0.954).
  \textbf{Hypothesis: plausible label noise} --- response
  shift/adaptation, stoic response style, or survey artifact. The model
  faithfully reports that this measured-health vector is overwhelmingly
  Fair/Poor-like; the disagreement is between the label and the
  respondent's own symptom answers.
\item
  \textbf{Case B --- most confident false negative} (row 392245, fold
  9): y = 1 (GENHLTH = 4 ``Fair'') yet P(y=1) = 0.0025 / 0.0027 /
  0.0156. A clean profile on every top feature: PHYSHLTH = 0 days (SHAP
  -0.61, the largest push toward 0), no chronic conditions, no walking
  difficulty, BMI 22.14, income bracket 9/11 (\$100--150k), college
  graduate, exercises. \emph{Every} top-12 SHAP feature contributes
  negatively (the uniformly blue middle panel of Figure 9 is the visual
  statement); $\Sigma$SHAP -3.37 $\rightarrow$ log-odds -5.98. This is systematic, not an
  oddball: among the 50 most confident shared errors, 50/50 have
  PHYSHLTH $\le$ 1 day, DIFFWALK = 2, DIABETE4 = 3, and 47/50 are GENHLTH =
  4. \textbf{Hypothesis: the error floor of the feature set} --- the
  ``Fair'' rating is anchored on constructs the survey does not measure
  (wellbeing beyond day counts, expectations, personality, social
  comparison); no model trained on these columns could rank this person
  highly without paying elsewhere.
\item
  \textbf{Case C --- the fence case} (row 208256, fold 4): y = 0
  (GENHLTH = 3) with LGBM P = 0.5000022 (XGB 0.598, LR 0.563);
  refit-proxy P = 0.481. A comorbidity block pushes toward y = 1
  (DIFFWALK SHAP +0.60, DIABETE4 +0.56, heart attack +0.36, stroke
  +0.30, ``unable to work'' +0.34, near-lowest income +0.23) while
  current-status anchors sit near zero (PHYSHLTH = 2 days, -0.02;
  POORHLTH = 0, -0.03); $\Sigma$SHAP +2.53, and the waterfall is visibly
  two-tailed. \textbf{Hypothesis: aleatoric} --- the population itself
  mixes 3s and 4s in this covariate region; P $\approx$ 0.5000 is the
  epistemically correct output, counted as an error only because
  thresholding must break exact ties.
\end{itemize}

\pandocbounded{\includegraphics[keepaspectratio]{C:/Users/Hassene/Documents/Machine learning final project/figures/p7_case_profiles.png}}

\emph{\textbf{Figure 9 --- Case profiles A/B/C.} Horizontal bars =
population z-scores of the case's raw values on the top-12
mean\textbar SHAP\textbar{} features; bar color = sign of that feature's
SHAP contribution (red pushes toward y=1, blue toward y=0); headers
carry GENHLTH and the three models' OOF probabilities. Source:
results/p7/selected\_cases.json, results/analysis/error\_analysis.md.}

\pandocbounded{\includegraphics[keepaspectratio]{C:/Users/Hassene/Documents/Machine learning final project/figures/p7_case_A_waterfall.png}}

\emph{\textbf{Figure 10 --- Case A local SHAP waterfall} (refit proxy):
PHYSHLTH +1.94 dominates the push from base -2.61 to log-odds +3.03.
Source: results/p7/shap\_cases.parquet.}

\pandocbounded{\includegraphics[keepaspectratio]{C:/Users/Hassene/Documents/Machine learning final project/figures/p7_case_B_waterfall.png}}

\emph{\textbf{Figure 11 --- Case B local SHAP waterfall}: uniformly
negative contributions --- nothing in the information set pushes toward
``Fair'', which }is* the explanation of the miss. Source:
results/p7/shap\_cases.parquet.*

\pandocbounded{\includegraphics[keepaspectratio]{C:/Users/Hassene/Documents/Machine learning final project/figures/p7_case_C_waterfall.png}}

\emph{\textbf{Figure 12 --- Case C local SHAP waterfall}: two-tailed ---
a comorbidity block versus near-zero current-status anchors, straddling
the 3\textbar4 fence. Source: results/p7/shap\_cases.parquet.}

\textbf{Global XAI (SHAP TreeExplainer, path-dependent, raw-margin
output; 20,000-row y-stratified sample, rs=42; additivity asserted to
8.8$\times$10\textsuperscript{-}\textsuperscript{1}\textsuperscript{4}; one-hot contributions summed back to the 76 original
columns).} Top-5 mean\textbar SHAP\textbar: \textbf{PHYSHLTH 0.6813,
DIFFWALK 0.2164, DIABETE4 0.1981, \texttt{\_BMI5} 0.1855, EDUCA 0.1786}
--- PHYSHLTH alone is 3.1$\times$ the \#2 feature (Figure 13).

Cross-method validation: SHAP ranks vs the P5 fold-internal gain ranks
agree at \textbf{Spearman $\rho$ = 0.911} (p = 3.6$\times$10\textsuperscript{-}\textsuperscript{3}\textsuperscript{0}, n = 76; Figure 14).
Two attribution mechanisms with different failure modes --- gain
double-counts correlated features; path-dependent SHAP spreads credit
along tree paths --- agree on the ordering. That agreement is strong
evidence that the driver ranking --- current physical-health status
first, functional limitation second, chronic-disease/SES context after
--- is a property of the data; it reproduces the EDA univariate ranking
(\S{}4) and the H1 blocks (\S{}7).

The model leans on self-reported current symptoms and therefore inherits
self-report bias --- a feature of GENHLTH as a target, not of LightGBM.

\pandocbounded{\includegraphics[keepaspectratio]{C:/Users/Hassene/Documents/Machine learning final project/figures/p7_shap_summary.png}}

\emph{\textbf{Figure 13 --- Global SHAP beeswarm}, top 20 original
features (contributions summed from one-hot children; log-odds output).
Source: results/p7/shap\_global\_importance.csv.}

\pandocbounded{\includegraphics[keepaspectratio]{C:/Users/Hassene/Documents/Machine learning final project/figures/p7_shap_vs_gain.png}}

\emph{\textbf{Figure 14 --- SHAP vs gain rank scatter}, Spearman $\rho$ =
0.911 (n = 76): two attribution methods agree on the ordering. Source:
results/p7\_summary.md, results/p5/lgbm\_importance.csv.}

Known caveat, stated rather than hidden: path-dependent TreeExplainer
spreads credit across correlated one-hot dummies, so per-case
attributions indicate \emph{which blocks} drive a prediction, not
precise per-feature credits (\texttt{results/p7/shap\_meta.json}); and
Case B's missing POORHLTH still contributes -0.12 through its NaN
one-hot category.

\section{Conclusions \& recommended
model}\label{conclusions-recommended-model}

\textbf{Recommended model: the tuned LightGBM pipeline of EXP-014a} ---
all-76 matrix v1 features, NaN-passthrough preprocessing,
\texttt{LGBMClassifier(n\_estimators=512,\ learning\_rate=0.0219,\ num\_leaves=94,\ min\_child\_samples=44,\ subsample=0.8159,\ colsample\_bytree=0.5980,\ reg\_lambda=3.123,\ subsample\_freq=1,\ random\_state=42)}.

Mean CV ROC-AUC \textbf{0.9043 $\pm$ 0.0017} on the canonical folds;
selection-bias-free nested estimate \textbf{0.9043 $\pm$ 0.0017} (EXP-015,
0.904303; the two agree to +0.000018). It is statistically significantly
better than the three alternatives it was formally tested against ---
its tuned XGBoost runner-up (T1), its own untuned configuration (T2),
and the tuned linear baseline (T3; Wilcoxon, all p = 0.001953, the n =
10 floor). Because the ten CV folds overlap, these paired p-values are
optimistic and are read here as strong evidence, not proof. The untested
families trail by clear margins --- RandomForest by +0.011 and the
single tree by +0.024, many times the fold noise --- while
HistGradientBoosting and untuned XGBoost sit within \textasciitilde0.001
and were not formally tested. It is also the fastest of the competitive
models (41.4 s per CV run untuned; 69.9 s tuned, vs RF's 351.5 s).

The recommendation rests on the full evidence chain, not one metric:
whitelist-controlled data (ADR-005) $\rightarrow$ measured 0.5 floor (EXP-001) $\rightarrow$
tuned linear/tree baselines (EXP-002/003) $\rightarrow$ six-family zoo with
identical folds (EXP-004..007) $\rightarrow$ imbalance verdict ``none'' with a
metric-level rationale (EXP-008..010) $\rightarrow$ selection/hypothesis/PCA all
answered honestly (EXP-011..013) $\rightarrow$ automated search + nested CV +
significance tests (EXP-014/015).

\textbf{What mattered, in order of effect size:}

\begin{enumerate}
\def\labelenumi{\arabic{enumi}.}
\tightlist
\item
  \emph{Family choice}: boosting over linear/bagging, +0.011 $\approx$ 6$\times$ fold
  std, significant (T3).
\item
  \emph{Data hygiene}: the whitelist removing a proven 100\%-agreement
  leak, plus the 88$\rightarrow$0 recode that flipped PHYSHLTH from $\rho$ -0.293 to
  +0.450 --- the single largest driver (51.69\% of gain).
\item
  \emph{Native missing-value routing}: NaN-passthrough kept module
  missingness as signal and is part of why boosters suit this survey.
\item
  \emph{Tuning}: real but small, +0.0010 (10/10 folds, T2) --- an order
  of magnitude below family choice.
\end{enumerate}

\textbf{Honest negative results, all reported with reasons:}

\begin{itemize}
\tightlist
\item
  Manual LogReg tuning was flat (the hyperparameter does not matter
  here).
\item
  SelectKBest(MI) lost to all-76 at every k (univariate filters drop
  correlated-but-useful columns trees use for free).
\item
  The H1-subset did \emph{not} strictly match all-76 within fold noise
  (99.2\% of lift, 0/10 folds --- kept all-76, used H1 for
  interpretation).
\item
  PCA cost 0.0146--0.0222 AUC (linear compression smears one-hot
  structure and forfeits NaN routing).
\item
  SMOTE actively hurt (-0.004) while class weighting and undersampling
  were AUC-neutral --- with the rank-based-metric explanation verified
  across two families.
\end{itemize}

\textbf{Limitations:}

\begin{itemize}
\tightlist
\item
  \emph{Boundary ambiguity}: 88.2\% of errors sit at the binarized
  3\textbar4 edge; a GENHLTH=4 rater with a near-identical covariate
  vector to a GENHLTH=3 rater plausibly exists (Case C), so the residual
  error is substantially aleatoric.
\item
  \emph{Self-report labels and features}: both target and top predictors
  are self-reported; the model inherits response shift, response style
  and recall bias (Case A).
\item
  \emph{Module missingness}: two of the five strongest univariate
  predictors (\texttt{JOINPAI2} 68.1\% missing, \texttt{POORHLTH}
  48.0\%) sit in high-missingness items, while the two strongest
  (\texttt{PHYSHLTH}, \texttt{DIFFWALK}) are near-complete core
  questions; imputed or NaN-routed, the high-missingness columns' signal
  is conditional on questionnaire design.
\item
  \emph{No causal claims}: this is associative prediction on a
  cross-sectional survey; a deployed system would additionally need
  calibration, subgroup fairness audits and a purpose-chosen threshold
  (Table 15 offers the menu).
\end{itemize}

\textbf{What would help:} modeling GENHLTH at the ordinal level
(e.g.~ordinal regression / ordinal-aware objectives) instead of
discarding the 1--5 ordering; richer mental-health and wellbeing
constructs (the error floor of Case B is exactly the missing construct);
and cost-sensitive thresholds tied to a real use case.

\textbf{Table 16 --- Comprehensive experiment summary.} All 42 rows of
\texttt{results/experiments\_log.csv} (mean CV ROC-AUC on the identical
canonical folds; SHA 3b3a17457039 asserted in all 41 model-experiment
rows --- EXP-000, the profiling row, has no folds by design).

{\def\LTcaptype{none} % do not increment counter
\begin{longtable}[]{@{}
  >{\raggedright\arraybackslash}p{(\linewidth - 6\tabcolsep) * \real{0.2500}}
  >{\raggedright\arraybackslash}p{(\linewidth - 6\tabcolsep) * \real{0.2500}}
  >{\raggedright\arraybackslash}p{(\linewidth - 6\tabcolsep) * \real{0.2500}}
  >{\raggedright\arraybackslash}p{(\linewidth - 6\tabcolsep) * \real{0.2500}}@{}}
\toprule\noalign{}
\begin{minipage}[b]{\linewidth}\raggedright
exp\_id
\end{minipage} & \begin{minipage}[b]{\linewidth}\raggedright
Model / run
\end{minipage} & \begin{minipage}[b]{\linewidth}\raggedright
Configuration / feature set
\end{minipage} & \begin{minipage}[b]{\linewidth}\raggedright
mean AUC
\end{minipage} \\
\midrule\noalign{}
\endhead
\bottomrule\noalign{}
\endlastfoot
EXP-000 & Data profiling & 438,693 $\times$ 303; leakage probe; no model &
--- \\
EXP-001 & DummyClassifier (prior) & all-76; sanity floor & 0.500000 \\
EXP-002a & LogisticRegression & C=0.01; all-76 & 0.892588 \\
EXP-002b & LogisticRegression & C=0.1 & 0.892580 \\
EXP-002c & LogisticRegression & C=1 & 0.892580 \\
EXP-002d & LogisticRegression & C=10 & 0.892580 \\
EXP-003a & DecisionTree & depth 6, leaf 1 & 0.848558 \\
EXP-003b & DecisionTree & depth 6, leaf 50 & 0.848645 \\
EXP-003c & DecisionTree & depth 6, leaf 200 & 0.848666 \\
EXP-003d & DecisionTree & depth 10, leaf 1 & 0.865805 \\
EXP-003e & DecisionTree & depth 10, leaf 50 & 0.871808 \\
EXP-003f & DecisionTree & depth 10, leaf 200 & 0.872352 \\
EXP-003g & DecisionTree & depth 16, leaf 1 & 0.795855 \\
EXP-003h & DecisionTree & depth 16, leaf 50 & 0.874830 \\
\textbf{EXP-003i} & DecisionTree & \textbf{depth 16, leaf 200 (winner)}
& \textbf{0.880206} \\
EXP-003j & DecisionTree & depth None, leaf 50 & 0.864202 \\
EXP-003k & DecisionTree & depth None, leaf 200 & 0.879195 \\
EXP-003l & DecisionTree & depth None, leaf 1 (defaults $\approx$ collapse) &
0.694922 \\
EXP-004 & RandomForest & 200 trees, leaf 50, sqrt & 0.893315 \\
EXP-005 & HistGradientBoosting & 300 iter, lr 0.06; NaN-passthrough &
0.903220 \\
EXP-006 & LightGBM & 500 trees, lr 0.05, 63 leaves & 0.903355 \\
EXP-007 & XGBoost & 500 trees, lr 0.05, depth 6 & 0.903357 \\
EXP-008a & XGBoost & + scale\_pos\_weight 5.0149 & 0.903261 \\
EXP-008b & LogisticRegression & + class\_weight balanced & 0.893035 \\
EXP-009a & XGBoost & + SMOTE(0.5) in folds & 0.899803 \\
EXP-009b & LogisticRegression & + SMOTE(0.5) in folds & 0.888419 \\
EXP-010a & XGBoost & + RandomUnderSampler(0.5) & 0.903096 \\
EXP-010b & LogisticRegression & + RandomUnderSampler(0.5) & 0.892765 \\
EXP-011a-50 & LightGBM & SelectKBest(MI) k=50 & 0.889178 \\
EXP-011a-100 & LightGBM & SelectKBest(MI) k=100 & 0.899556 \\
EXP-011a-200 & LightGBM & SelectKBest(MI) k=200 & 0.899670 \\
EXP-011b & LightGBM & all-76; fold-internal gain importances &
0.903355 \\
EXP-012a & LightGBM & H1-subset (36 cols) & 0.900051 \\
EXP-012b & XGBoost & H1-subset (36 cols) & 0.900316 \\
EXP-012c & LightGBM & complement (40 cols) & 0.811844 \\
EXP-012d & XGBoost & complement (40 cols) & 0.811881 \\
EXP-013-50 & LightGBM & PCA 50 (43.7\% var) & 0.881126 \\
EXP-013-100 & LightGBM & PCA 100 (63.7\% var) & 0.886203 \\
EXP-013-200 & LightGBM & PCA 200 (97.8\% var) & 0.888800 \\
\textbf{EXP-014a} & \textbf{LightGBM (tuned)} &
\textbf{RandomizedSearchCV winner --- final model} &
\textbf{0.904321} \\
EXP-014b & XGBoost (tuned) & RandomizedSearchCV winner & 0.903843 \\
EXP-015 & LightGBM nested CV & outer = canonical folds; unbiased tuned
estimate & 0.904303 \\
\end{longtable}
}

\section{Ethics and Use of Generative
AI}\label{ethics-and-use-of-generative-ai}

This section is the mandatory final disclosure, drafted from the
project's running AI-assistance log (23 entries, 2026-10-25 to
2026-11-08) and following the disclosure structure required by the
course guidelines: tool used, purpose, and how the generated output was
used in the final work.

\textbf{Statement of authorship.} Generative AI was used substantially
and systematically throughout this project, and its use is disclosed
here in full. All AI-generated contributions --- plans, code, analysis,
and prose --- were reviewed by the student; methodological decisions
were ratified through the project's decision-record process (ADRs
001--005), and every experimental number in this report is
machine-traceable to \texttt{results/experiments\_log.csv} or to a named
committed artifact, independently of any AI-generated text.
\textbf{Hassene Bouamama assumes responsibility for and authorship of
the submitted work.}

\textbf{Assistance from other persons.} No assistance beyond course
instruction was received.

\textbf{Tool used.} The sole generative-AI tool was an \textbf{OpenCode
coding agent} (model: GLM, by Z.ai), operated through a
Claude-Code-style agentic harness inside the project repository. The
agent worked in narrowly scoped roles (planning, implementation,
analysis, reporting) under a written rulebook (\texttt{AGENTS.md}) that
fixed the target, folds, seeds, metric and logging contract before any
model was run; the student directed, reviewed and ratified its outputs
at each milestone. No other generative-AI tools (ChatGPT, Copilot,
Gemini, etc.) were used.

\textbf{Full disclosure by log entry (all 23 entries, verbatim in
substance):}

{\def\LTcaptype{none} % do not increment counter
\begin{longtable}[]{@{}
  >{\raggedright\arraybackslash}p{(\linewidth - 8\tabcolsep) * \real{0.0612}}
  >{\raggedright\arraybackslash}p{(\linewidth - 8\tabcolsep) * \real{0.1224}}
  >{\raggedright\arraybackslash}p{(\linewidth - 8\tabcolsep) * \real{0.1224}}
  >{\raggedright\arraybackslash}p{(\linewidth - 8\tabcolsep) * \real{0.1837}}
  >{\raggedright\arraybackslash}p{(\linewidth - 8\tabcolsep) * \real{0.5102}}@{}}
\toprule\noalign{}
\begin{minipage}[b]{\linewidth}\raggedright
\#
\end{minipage} & \begin{minipage}[b]{\linewidth}\raggedright
Date
\end{minipage} & \begin{minipage}[b]{\linewidth}\raggedright
Tool
\end{minipage} & \begin{minipage}[b]{\linewidth}\raggedright
Purpose
\end{minipage} & \begin{minipage}[b]{\linewidth}\raggedright
How the output was used
\end{minipage} \\
\midrule\noalign{}
\endhead
\bottomrule\noalign{}
\endlastfoot
1 & 2026-10-25 & OpenCode agent (GLM, Z.ai) & Project planning: read
guidelines + codebook, framed binary target, designed pipeline, agent
team, repo and vault structure & Plan reviewed and approved by the
student; structure adopted for the project \\
2 & 2026-10-25 & OpenCode agent (GLM, Z.ai) & Repository scaffolding:
.gitignore, README, AGENTS.md rules, vault seed content, src/config.py &
Files reviewed and committed as the initial project setup \\
3 & 2026-10-25 & OpenCode agent (GLM, Z.ai) & Interpreting the BRFSS
codebook structure (missing-value codes, section layout, GENHLTH
distribution) & Used to draft the preprocessing rules and leakage
watchlist; verified against the data in EXP-000 \\
4 & 2026-10-25 & OpenCode agent (GLM, Z.ai) & EXP-000 profiling script +
analysis of missingness/association/leakage probe & Script committed as
src/profile\_data.py; findings documented in vault and verified in the
P2 matrix build \\
5 & 2026-10-25 & OpenCode agent (GLM, Z.ai) & P2 whitelist construction
+ preprocessing implementation (src/preprocessing.py) +
sentinel/codebook cross-verification & Code reviewed via sanity gates;
all 303 column verdicts in results/profiling/triage\_decisions.csv;
decisions ratified as ADR-005 \\
6 & 2026-10-25 & OpenCode agent (GLM, Z.ai) & EXP-001..003
implementation + run via src/experiments.py; fold-hash verification;
manual hyperparameter grids & Used as the committed experiment harness
(resume-safe, fold-sha asserted) with results logged in
results/experiments\_log.csv; findings documented in vault EXP notes \\
7 & 2026-10-25 & OpenCode agent (GLM, Z.ai) & EDA figures +
association/correlation analysis (results/analysis/eda\_notes.md) &
Figures committed to figures/; findings cited in this report after human
visual review \\
8 & 2026-11-02 & OpenCode agent (GLM, Z.ai) & P4 model zoo
implementation + runs (src/p4\_models.py), imbalance comparison,
NaN-passthrough verification, harness parser fix & Results logged in
experiments\_log.csv; verdicts documented in vault; cited in report
\S{}6 \\
9 & 2026-11-02 & OpenCode agent (GLM, Z.ai) & P5 feature-selection
experiments (SelectKBest, fold-internal gain importances, H1/complement
subsets, PCA), figures and summary & Results logged in
experiments\_log.csv; H1 verdict + importance findings documented in
vault for report citation (\S{}7) \\
10 & 2026-11-02 & OpenCode agent (GLM, Z.ai) & P6 randomized searches,
tuned finals, nested CV, Wilcoxon tests (src/p6\_tune.py) & Results
logged; statistical verdicts documented for report (\S{}8) \\
11 & 2026-11-02 & OpenCode agent (GLM, Z.ai) & P7a computation: OOF
prediction pipeline for 3 logged configs (per-fold AUC re-verified
$\le$4.8e-7), error taxonomy + cross-model wrongness, case selection, SHAP
TreeExplainer global/local with dummy-group aggregation
(src/p7\_error\_xai.py) & Artifacts committed (results/p7/,
figures/p7\_*); facts documented in vault; cited in report \S{}9 \\
12 & 2026-11-02 & OpenCode agent (GLM, Z.ai) & P7b interpretation:
case-profile figure, 50-row subtype profiling, threshold sweep,
error-analysis note with FACT/INTERPRETATION discipline
(src/analysis/p7b\_cases.py, results/analysis/error\_analysis.md) &
Hypotheses (label noise / error floor / aleatoric boundary) documented
as interpretations for report drafting; figures pending human visual
pass \\
13 & 2026-11-02 & OpenCode agent (GLM, Z.ai) & P8a deliverable assembly:
deliverables/code.py (self-contained single-file pipeline) +
verification runs (--smoke, --full twice; per-fold deltas vs EXP-014a $\le$
4.84e-7, folds-sha asserted) & Committed as the required end-to-end
deliverable; verification numbers cited in report \S{}10 \\
14 & 2026-11-02 & OpenCode agent (GLM, Z.ai) & P8b full report drafting
(11 sections, 16 tables, 14 figures, \S{}10 final-model justification, \S{}11
Ethics from this log), PDF build chain, and the 6-beat screencast script
& Committed as deliverables/report/report.pdf v1.0 +
screencast\_script.md; \textasciitilde45 quantitative claims QA-traced
to the experiment log with 0 mismatches \\
15 & 2026-11-02 & OpenCode agent (GLM, Z.ai) & QA remediation edits:
fold-SHA coverage wording corrected in two report places + PDF rebuild &
Verified in the rebuilt PDF text layer; closes QA must-fix \#2 \\
16 & 2026-11-02 & OpenCode agent (GLM, Z.ai) & Pre-submission QA audit
(read-only): rubric-evidence verification, independent sandbox
reproduction of code.py --full fold AUCs ($\le$ 4.84e-7) and the confusion
matrix, leakage/traceability/git-history checks & Audit verdict OVERALL
PASS accepted as the pre-submission gate; findings recorded and
resolved/deferred in the P8 daily log \\
17 & 2026-11-03 & OpenCode agent (GLM, Z.ai) & Personal learning
notebook notebooks/BRFSS\_walkthrough.ipynb: 14-section full-data replay
of the pipeline (canonical folds and src/ pipelines reused, never
re-implemented), executed 41.3 min with 0 cell errors, 12/12 self-check
rows matching the experiment log; the verification pass uncovered a
correlation-claim scope error & Committed as a personal learning
artifact (header-marked, not a deliverable); the scope discovery
triggered the entry-18 correction \\
18 & 2026-11-03 & OpenCode agent (GLM, Z.ai) & Report correlation-claim
remediation: \S{}4 and Figure-4 caption re-scoped (top-40 view vs full
76$\times$76 maximum) + PDF rebuild & Committed; text-layer diff verified to
contain only the two intended edits document-wide \\
19 & 2026-11-08 & OpenCode agent (GLM, Z.ai) & Submission-day polish of
deliverables/code.py: comment-only restyle (compact header, concise
section banners); token-skeleton verification proved zero executable
changes; --smoke run re-passed end-to-end & Restyled code.py reviewed by
the student and verified before submission \\
20 & 2026-11-08 & OpenCode agent (GLM, Z.ai) & Submission-day airing of
this report: reformatted \S{}1--\S{}11 into short paragraphs/bullets with one
citation per claim; author name filled; numeric freeze machine-verified
(all numeric tokens, EXP ids, tables and figures unchanged) &
Reformatted report reviewed by the student; compiled to report.pdf \\
21 & 2026-11-08 & OpenCode agent (GLM, Z.ai) & Screencast preparation:
9-slide deck (screencast\_deck.pptx) from the logged results with
speaker notes; rehearsal aid mapping each claim to the student-run
verification & Student rehearsed with the deck and recorded the
screencast \\
22 & 2026-11-08 & OpenCode agent (GLM, Z.ai) & Pre-submission
claim-audit fixes from the student's critical review (12 items):
statistical-scope wording, SMOTE/sklearn/booster overclaims,
univariate-missingness claim, RandomForest-NaN justification (sklearn
1.9), POORHLTH reconciliation, ``bimodal'' softened, ALCDAY5 relabelled,
H1 block wording, nested-CV budget caveat, family counting & Corrections
applied to report.md and code.py labels; PDF rebuilt \\
23 & 2026-11-08 & OpenCode agent (GLM, Z.ai) & EDA stage added to
code.py (target imbalance, missingness profile, top associations;
regenerates the three core EDA figures in full mode) & Prompted by the
student's pre-submission review of the deliverable code; smoke and full
runs verified \\
\end{longtable}
}

\textbf{How AI assistance was bounded.} The rulebook made the
\emph{scientific} decisions human-ratified (ADRs) rather than
agent-made. The append-only experiment log and per-run fold-hash
assertions mean the report's numbers stand on the recorded runs, not on
the agent's word. Interpretations from AI-assisted analysis are labeled
as hypotheses in the project notes (FACT/INTERPRETATION discipline in
the error-analysis note) and are presented as such in \S{}9. This report
itself was drafted by the agent role from the log and artifacts, then is
to be reviewed and personalized by the student before submission.

\begin{center}\rule{0.5\linewidth}{0.5pt}\end{center}

\emph{End of report. Build:
\texttt{deliverables/report/build\_report.py} $\rightarrow$ \texttt{report.pdf}.
Sources: \texttt{results/experiments\_log.csv} (42 rows),
\texttt{results/p3\_baselines\_summary.md},
\texttt{results/p4\_summary.md}, \texttt{results/p5\_summary.md},
\texttt{results/p6\_summary.md}, \texttt{results/p7\_summary.md},
\texttt{results/analysis/eda\_notes.md},
\texttt{results/analysis/error\_analysis.md},
\texttt{results/p6/wilcoxon\_tests.csv}, \texttt{results/p7/*}, vault
notes (ADRs 001--005, EXP notes, data notes).}
\end{document}